\SourceLeaf{232}%
Rechnung einer neuen Ephemeride ganz danach zu wiederholen. Es wird alsdann hinreichend sein, \textit{einige} Örter des Planeten nachzurechnen; denn der Unterschied kann nicht sehr beträchtlich sein, und nur in langsamen Stufen regelmässig anwachsen. So hat schon Dr.\ \textsc{Gauss} es versucht, einige Orte der \textit{Pallas} nach seinen III. Elementen zu berechnen, und der Unterschied ist eben nicht so beträchtlich, dass dadurch die Auffindung des Planeten gehindert werden sollte, wie man aus beikommender Vergleichung ersehen kann, [welcher hier die Resultate für die Elemente V.\ aus der Handschrift beigefügt sind.]
\begin{center}
\begingroup\small\setlength{\tabcolsep}{2.5pt}
\begin{tabular}{l|r|r|r|r|r|r|r|r}
&\multicolumn{2}{c|}{AR.\ der $\PallasSymbol$}&&\multicolumn{2}{c|}{Abweich.\ der $\PallasSymbol$}&&AR.&Decl.\\
1803&III. Elem.&IV. Elem.&Diff.&III. Elem.&IV. Elem.&Diff.&V. Elem.&V. Elem.\\\hline
Febr.\ 4&$267^0\ 35'$&$267^0\ 48'$&$13'$&$5^0\ 41'$&$5^0\ 38'$&$3'$&$267^0\ 45'$&$5^0\ 39'$\\
Juni 28&$275\ 45$&$276\ 17$&$32$&$23\ 12$&$23\ 10$&$2$&$276\ 9$&$23\ 11$
\end{tabular}\endgroup
\end{center}
Der Gegenschein dieses Planeten fällt nach den III. Elementen den 30. Juni bürgerl.\ Zeit Vormittags, nach den IV. Elementen an eben dem Tage Nachmittags. Der Unterschied in der Abweichung ist ganz unbedeutend. Nach V.\ Juni 30.\ $2^{\mathrm h}\ 53^{\mathrm m}$ In $277^0\ 45'$ L.
\begin{flushright}[\textsc{von Zach.}]\end{flushright}

\bigskip\Needspace{10\baselineskip}
\pdfbookmark[1]{Ceres Ferdinandea. 1802 November}{ceres180211}
\begin{center}
CERES FERDINANDEA.\par
\medskip\hrule\medskip
Monatliche Correspondenz. Band VI. S.\ 492—498. 1802 November.\par
\medskip\hrule
\end{center}\medskip
—\quad —\quad —

Dr.\ \textsc{Gauss} hat indessen fortgesetzt an der Bahn dieses Planeten gearbeitet. Mit Zuziehung der Störungen, so wie er sie (Oct.\ St.\ S.\ 387) [224 d.\ B.] nach seinen VII. Elementen berechnet hatte, in Verbindung mit den Breiten-Gleichungen, bei denen seine Rechnung mit der \textsc{Oriani}'schen (Jun.\ St.\ S.\ 586. Jul.\ St.\ S.\ 68) [1802. Juli S.\ 67—70. —\quad —\quad —\quad \textsc{Oriani} dieser geschickte, und in den verwickelten Störungsrechnungen so sehr gewandte Calculator, liess sich die Mühe nicht verdriessen, die ganze Perturbationsrechnung der \textit{Ceres} in einer andern Hypothese der mittlern Entfernung zu wiederholen, und dadurch zugleich seine erst erhaltenen Resultate nochmals zu prüfen. Hier ist das Resultat dieser ganzen geführten Rechnung, welche uns dieser grosse Astronom mitzutheilen die Güte hatte.
\[
\begin{aligned}
\text{Es sei }D&=\text{mittlere Länge der }\CeresSymbol-\text{mittlere Länge des }\JupiterSymbol;\\
A'&=\text{mittl. Anomalie }\JupiterSymbol;\\
A&=\text{mittl. Anomalie der }\CeresSymbol.\\
H'&=\text{mittl. Länge }\JupiterSymbol-\text{mittl. Länge }\AscendingNode\JupiterSymbol\\
H&=\text{mittl. Länge }\CeresSymbol-\text{mittl. Länge }\AscendingNode\CeresSymbol.
\end{aligned}
\]
Ferner ist zu bemerken, dass man die Störungs-Gleichungen für jede andere Excentricität $e$ der Bahn erhält, wenn man die Glieder, welche $A$ enthalten, multiplicirt mit $\dfrac{e}{0.081406}$; die Glieder, welche $2A$ enthalten, multiplicirt mit $\left(\dfrac{e}{0.081406}\right)^2$; die Glieder, welche $3A$ enthalten. multiplicirt mit $\left(\dfrac{e}{0.081406}\right)^3$.

\Needspace{24\baselineskip}
Die Störungs-Gleichungen der \textit{Ceres} durch Jupiter und ihre Änderungen sind alsdann wie folgt:
\SourceLeaf{233}%
Die Störungen der heliocentrischen Breite sind:
\begin{center}
\begingroup\small
\begin{tabular}{l|r}
\shortstack{Nach den VII. \textsc{Gauss}'schen\\Elementen der Bahn}&\shortstack{Wenn die mittlere jährliche\\Bewegung der Ceres um\\20 Min. vermehrt wird.}\\\hline
$-13''19\sin(H-D)$&$-12''97$\\
$+1.62\sin H'$&$+1.59$\\[3pt]
$+16.21\sin(2D-H)$&$+16.03$\\
$-1.99\sin(D-H')$&$-1.96$\\[3pt]
$+32.38\sin(3D-H)$&$+32.78$\\
$-3.97\sin(2D-H')$&$-4.01$\\[3pt]
$-5.50\sin(4D-H)$&$-5.31$\\
$+0.67\sin(3D-H')$&$+0.65$\\[3pt]
$+6.17\sin(D+H)$&$+6.06$\\
$-0.76\sin(2D+H')$&$-0.74$\\[3pt]
$+1.52\sin(2D+H)$&$+1.49$\\
$-0.19\sin(5D+H')$&$-0.18$\\[3pt]
$+16.15\sin(4D-A'-H-2^0\ 33')$&$+13.83$\\
$-13.86\sin(5D-A-H-2^0\ 33')$&$-11.92$
\end{tabular}\endgroup
\end{center}
fast vollkommen harmonirt] bestimmte er sogleich folgende Elemente:
\begin{center}
\begingroup\small\setlength{\tabcolsep}{3pt}
\begin{tabular}{l r}
Epoche 1801 für Seeberg $\ .\ .$&$77^0\ 19'\ 38''4$\\
tägliche tropische Bewegung $\ .$&$770''764$\\
Umlaufszeit $\ .\ .\ .\ .\ .\ .\ .$&1681 Tage 11 Stund.\\
Excentricität $\ .\ .\ .\ .\ .\ .\ .$&$0.0788132$\\[3pt]
Logarithmus der halben grossen Axe $\ .$&$0.4421085$\\
Aphelium 1801 $\ .\ .\ .\ .\ .\ .\ .\ .\ .$&$326^0\ 33'\ 10''$\\
Knoten 1801 $\ .\ .\ .\ .\ .\ .\ .\ .\ .\ .$&$80\ 54\ 52$\\
Neigung — $\ .\ .\ .\ .\ .\ .\ .\ .\ .\ .$&$10\ 37\ 48$
\end{tabular}\endgroup
\end{center}

Mit diesen Elementen berechnete nun Dr.\ \textsc{Gauss} die Störungen der $\CeresSymbol$ durch $\JupiterSymbol$ aufs neue, wie folgt.
\begin{center}\textit{Secular-Gleichungen.}\end{center}
\begin{center}
\begingroup\small\setlength{\tabcolsep}{3pt}
\begin{tabular}{l r}
Jährl. Vorrücken der Sonnenferne gegen die Nachtgleichen&$+121''3$\\
Jährl. Vorrücken der Knoten gegen die Nachtgleichen $\ .\ .$&$+0.8$\\
Jährl. Abnahme der Excentricität $\ .\ .\ .\ .\ .\ .\ .\ .\ .$&$0.000005829$\\
Jährl. Abnahme der Neigung der Bahn $\ .\ .\ .\ .\ .\ .\ .$&$+0''38$
\end{tabular}\endgroup
\end{center}
\Needspace{19\baselineskip}
\begin{center}\textit{Periodische Gleichungen.}\par\medskip
\textbf{In der Länge:}\end{center}
\[
\begin{aligned}
&-231''00\sin(\CeresSymbol-\JupiterSymbol)\\
&+492.37\sin2(\CeresSymbol-\JupiterSymbol)\\
&+43.85\sin3(\CeresSymbol-\JupiterSymbol)\\
&+10.00\sin4(\CeresSymbol-\JupiterSymbol)\\
&+3.03\sin5(\CeresSymbol-\JupiterSymbol)\\
&+1.04\sin6(\CeresSymbol-\JupiterSymbol)\\
&+0.40\sin7(\CeresSymbol-\JupiterSymbol)\\
&+0.16\sin8(\CeresSymbol-\JupiterSymbol)\\
&-58.90\sin(\JupiterSymbol+16^0\ 56'\ 37'')\\
&-598.69\sin(\CeresSymbol-2\JupiterSymbol-26^0\ 6'\ 58'')\\
&-437.75\sin(2\CeresSymbol-3\JupiterSymbol-10^0\ 54'\ 25'')
\end{aligned}
\]
\[
\begin{aligned}
&+54''90\sin(3\CeresSymbol-4\JupiterSymbol-10^0\ 32'\ 23'')\\
&+10.28\sin(4\CeresSymbol-5\JupiterSymbol-10\quad21\quad35'')\\
&+3.19\sin(5\CeresSymbol-6\JupiterSymbol-10\quad15\quad2'')\\
&+1.20\sin(6\CeresSymbol-7\JupiterSymbol-10\quad10\quad20'')\\
&+0.50\sin(7\CeresSymbol-8\JupiterSymbol-10\quad6\quad24'')\\
&+22.72\sin(2\CeresSymbol-\JupiterSymbol+36\quad0\quad12'')\\
&-51.80\sin(3\CeresSymbol-2\JupiterSymbol+32\quad47\quad57'')\\
&-5.74\sin(4\CeresSymbol-3\JupiterSymbol+30\quad41\quad32'')\\
&-1.65\sin(5\CeresSymbol-4\JupiterSymbol+28\quad53\quad58'')\\
&-0.69\sin(6\CeresSymbol-5\JupiterSymbol+28\quad13\quad52'')\\
&-0.26\sin(7\CeresSymbol-6\JupiterSymbol+26\quad41\quad16'')
\end{aligned}
\]
\SourceLeaf{234}%
\Needspace{19\baselineskip}
\begin{center}\textbf{Im Radius vector:}\end{center}
\[
\begin{aligned}
&-0.0000943\\
&+0.0010261\cos(\CeresSymbol-\JupiterSymbol)\\
&-0.0037679\cos2(\CeresSymbol-\JupiterSymbol)\\
&-0.0004176\cos3(\CeresSymbol-\JupiterSymbol)\\
&-0.0001069\cos4(\CeresSymbol-\JupiterSymbol)\\
&-0.0000347\cos5(\CeresSymbol-\JupiterSymbol)\\
&-0.0000127\cos6(\CeresSymbol-\JupiterSymbol)\\
&-0.0000050\cos7(\CeresSymbol-\JupiterSymbol)\\
&-0.0000021\cos8(\CeresSymbol-\JupiterSymbol)\\
&-0.0000614\cos(\CeresSymbol+24^0\ 13'\ 37'')\\
&+0.0002398\cos(\JupiterSymbol+22\quad58\quad33'')\\
&+0.0008325\cos(\CeresSymbol-2\JupiterSymbol-23^0\ 54'\ 51'')
\end{aligned}
\]
\[
\begin{aligned}
&+0.0025560\cos(2\CeresSymbol-3\JupiterSymbol-10^0\ 39'\ 20'')\\
&-0.0004698\cos(3\CeresSymbol-4\JupiterSymbol-10\quad51\quad24'')\\
&-0.0001048\cos(4\CeresSymbol-5\JupiterSymbol-10\quad46\quad51'')\\
&-0.0000358\cos(5\CeresSymbol-6\JupiterSymbol-10\quad40\quad14'')\\
&-0.0000142\cos(6\CeresSymbol-7\JupiterSymbol-10\quad34\quad4'')\\
&-0.0000061\cos(7\CeresSymbol-8\JupiterSymbol-10\quad28\quad56'')\\
&-0.0001313\cos(2\CeresSymbol-\JupiterSymbol+37\quad11\quad18'')\\
&+0.0002928\cos(3\CeresSymbol-2\JupiterSymbol+32\quad17\quad53'')\\
&+0.0000397\cos(4\CeresSymbol-3\JupiterSymbol+29\quad3\quad25'')\\
&+0.0000135\cos(5\CeresSymbol-4\JupiterSymbol+26\quad58\quad26'')\\
&+0.0000057\cos(6\CeresSymbol-5\JupiterSymbol+25\quad46\quad37'')\\
&+0.0000026\cos(7\CeresSymbol-6\JupiterSymbol+25\quad2\quad49'')
\end{aligned}
\]
\begin{center}\textbf{In der Breite.}\end{center}
\[
\begin{aligned}
&-11''60\sin(\JupiterSymbol-78^0\ 31'\ 50'')\\
&+14.29\sin(\CeresSymbol-2\JupiterSymbol+78^0\ 31'\ 50'')\\
&+28.73\sin(2\CeresSymbol-3\JupiterSymbol+78\quad31\quad50'')\\
&-4.81\sin(3\CeresSymbol-4\JupiterSymbol+78\quad31\quad50'')\\
&-0.96\sin(4\CeresSymbol-5\JupiterSymbol+78\quad31\quad50'')
\end{aligned}
\]
\[
\begin{aligned}
&-0''29\sin(5\CeresSymbol-6\JupiterSymbol+78^0\ 31'\ 50'')\\
&+5.43\sin(2\CeresSymbol-\JupiterSymbol-78\quad31\quad50'')\\
&+1.34\sin(3\CeresSymbol-2\JupiterSymbol-78\quad31\quad50'')\\
&+0.44\sin(4\CeresSymbol-3\JupiterSymbol-78\quad31\quad50'')
\end{aligned}
\]
Endlich berechnete Dr.\ \textsc{Gauss} nach diesen Formeln den numerischen Werth der Störungen für die zum Grunde gelegten Beobachtungen aufs neue, und fand damit folgende Elemente der Bahn, die wir künftig mit (VIII) bezeichnen werden.
\begin{center}
\begin{tabular}{r|r|r|r}
\multicolumn{2}{c|}{Epoche für Seeberg}&Sonnenferne&Knoten\\\hline
1801&$77^0\ 19'\ 34''9$&$326^0\ 33'\ 37''$&$80^0\ 54'\ 59''$\\
1802&$155\ 28\ 35.1$&$326\ 35\ 39$&$80\ 55\ 0$\\
1803&$233\ 37\ 35.3$&$326\ 37\ 40$&$80\ 55\ 1$
\end{tabular}
\end{center}
\begin{center}
\begingroup\small\setlength{\tabcolsep}{3pt}
\begin{tabular}{l r}
tägliche tropische Bewegung $\ .\ .\ .\ .\ .\ .\ .\ .\ .\ .\ .$&$770''7951$\\
Tropische Umlaufszeit $\ .\ .\ .\ .\ .\ .\ .\ .\ .\ .\ .\ .\ .$&1681 Tage 9 Stunden\\
Excentricität (1801) $\ .\ .\ .\ .\ .\ .\ .\ .\ .\ .\ .\ .\ .\ .$&$0.0788352$\\
Neigung (1801) $\ .\ .\ .\ .\ .\ .\ .\ .\ .\ .\ .\ .\ .\ .\ .\ .$&$10^0\ 37'\ 56''0$\\
Logarithmus der halben Axe $\ .\ .\ .\ .\ .\ .\ .\ .\ .\ .$&$0.4420971$
\end{tabular}\endgroup
\end{center}
Diese neuen Elemente sind von den vorhergehenden so sehr wenig verschieden, dass man eine neue Berechnung der Störungen nach denselben als überflüssig ansehen kann, da sie von den obigen nicht merklich verschieden ausfallen können. Aus Neugierde hat Dr.\ \textsc{Gauss} es versucht, wie viel diese neuen VIII. Elemente von den VII. gegen die Zeit der Zusammenkunft im künftigen Jahre abweichen werden. Er findet mit Übergehung der Aberration für Juni 28, 1803 um 12 Uhr mittl.\ Seeberger Zeit
\begin{center}
\begin{tabular}{c|r|r}
\shortstack{Nach den\\Elementen}&Länge der $\CeresSymbol$&\shortstack{Breite der $\CeresSymbol$\\südl.}\\\hline
VII&$280^0\ 7'\ 43''$&$5^0\ 1'\ 21''$\\
VIII&$280\ 17\ 58$&$5\ 4\ 53$\\\hline
Unterschied&$10\ 15$&$3\ 32$
\end{tabular}
\end{center}
\expandafter\def\csname HistoricalNoteText22\endcsname{Observations on the two lately discovered celestial Bodies, by \textsc{William Herschel}. Read before the Royal Society. May 6. 1802. Aus den Philos.\ Transact.\ MDCCCII.}
\SourceLeaf{235}%

Dr.\ \textsc{Gauss} schliesst seine Berechnung mit folgender Bemerkung: „Bei der grossen Mühe, die die Berechnung eines Ortes des Planeten schon nach obigen Störungs-Gleichungen macht, scheint es wol noch eben nicht rathsam, ihre Anzahl schon jetzt durch die Saturns-Gleichungen und durch die von den Quadraten und Producten der Excentricitäten abhängigen beim Jupiter zu vermehren, deren Aufnahme übrigens weiter keine Schwierigkeit, als das \textit{Taedium} des mechanischen Calculs haben würde, zumal da wir nun beinahe nach einem halben Jahre mit Hülfe neuer Beobachtungen den Elementen einen viel grösseren Grad von Schärfe werden geben können. Da die Berechnung des jedesmaligen Betrags der Gleichungen so sehr beschwerlich ist: so wird es, wie mich dünkt, von Wichtigkeit sein, auf Abkürzung und Vereinfachung der Tafeln für dieselben zu denken, worüber ich Ihnen meine Ideen künftig nebst einem \textit{Specimen} vorlegen werde.“
\begin{flushright}[\textsc{von Zach.}]\end{flushright}

\bigskip\Needspace{10\baselineskip}
\pdfbookmark[1]{Pallas Olbersiana. 1802 November}{pallas180211}
\begin{center}
PALLAS OLBERSIANA.\par
\medskip\hrule\medskip
Monatliche Correspondenz. Band VI. S.\ 499—505. 1802 November.\par
\medskip\hrule
\end{center}\medskip
—\quad —\quad —

Dr.\ \textsc{Herschel} hat die Güte gehabt, uns einen besondern Abdruck seiner Abhandlung aus den \textit{Philosophischen Transactionen}\HistoricalNote{22} zuzuschicken, deren wesentlichen Inhalt wir bereits unseren Lesern im Auszug aus einem seiner Briefe im Juli-Stück S.\ 90 mitgetheilt haben. Nach nochmaliger aufmerksamer Durchlesung dieser Abhandlung finden wir uns nicht bewogen, unsere darüber schon geäusserte Meinung zu ändern; auch sehen wir aus unserm Briefwechsel, dass auch andere Astronomen, wie z.\ B.\ \textsc{La Lande}, \textsc{Oriani}, \textsc{Piazzi}, \textsc{Gauss} u.\ a.\ m.\ dieser seiner neuen vorgeschlagenen Classification nicht beipflichten. Besonders verdienen die Gründe und eine recht treffliche Bemerkung des Dr.\ \textsc{Gauss} gegen das angebliche sogenannte harmonische Gesetz, womit die Abstände der Planeten übereinstimmen sollen, wohl erwogen zu werden. Dieser scharfsinnige Messkünstler drückt sich hierüber in einem Schreiben vom 16. October auf folgende merkwürdige Art aus.

„Dr.\ \textsc{Herschel} will, wie mir Prof.\ \textsc{Huth} sagte, der ihn in England besucht hat, noch immer die neuen Planeten nicht toleriren, obgleich meines Wissens noch nicht ein einziger Astronom seinen Vorschlag gebilligt hat. Im Grunde hängt es, wie Sie so richtig bemerkt haben, nur von der Übereinkunft ab, ob wir die \textit{Ceres} und die \textit{Pallas} als Planeten ansehen wollen oder nicht, und es ist gar nicht die Rede davon, ob sie Planeten \textit{sind} oder nicht, sondern ob es \textit{schicklich} und \textit{passlich} ist, diese Weltkörper, die in einigen Punkten mit den \textit{bisher bekannten} Planeten übereinstimmen, in andern von ihnen abgehen, Planeten zu \textit{nennen}, oder nicht. Dass nun die letzteren Punkte (unter welchen bei der \textit{Ceres} die Neigung der Bahn gar nicht einmal gerechnet werden kann, wie Sie im December-Hefte des vorigen Jahres bemerkt haben) ganz unwesentlich sind, haben Sie zur Genüge im Juli-Heft dieses Jahrs gezeigt, und dass bei den Astronomen eine kreisähnliche Bahn, und die davon abhängige \textit{perennirende Gegenwart} immer als das Wesentliche gegolten habe, scheint der Umstand zu beweisen, dass
\expandafter\def\csname HistoricalNoteText23\endcsname{[Vergl.\ Astronomisches Jahrbuch für 1790. \textsc{Wurm} an \textsc{Bode}, Leonberg 1787 Februar 27. Von möglichen Planeten und Cometen unsers Sonnensystems. S.\ 168.]}
\SourceLeaf{236}%
die Astronomen sogleich den Planetismus ohne weiteres anerkannt haben, sobald sie sich von jener Beschaffenheit der Bahn überzeugt hielten. Mich dünkt sogar, dass wir, wenn die Zukunft die Vermuthung unseres vortrefflichen \textsc{Olbers} \textit{Ceres} und \textit{Pallas} seien nur Stücke von einem zerstörten Planeten, durch die Auffindung anderer Stücke bestätigen sollte, selbst in diesem Falle noch eben nicht nothwendig von dem Namen Planeten abgehen müssen; es scheint mir nemlich bei der Untersuchung der Ansprüche auf den Titel eines Planeten mehr darauf anzukommen, \textit{ob} diese Weltkörper nach ihrer wesentlichen Eigenschaft Planeten sind, als \textit{wie} sie es \textit{geworden} sind.

„Sonderbar ist es, dass man das vom Prof.\ \textsc{Titius} angegebene sogenannte Gesetz als ein Argument gegen die beiden Planeten gebrauchen wollte. Dieses Verhältniss trifft bei den übrigen Planeten gegen die Natur aller Wahrheiten, die den Namen Gesetze verdienen, nur ganz beiläufig, und, was man noch nicht einmal bemerkt zu haben scheint\HistoricalNote{23}, beim Mercur gar nicht zu. Es scheint mir sehr einleuchtend, dass die Reihe
\[
4,\ 4+3,\ 4+6,\ 4+12,\ 4+24,\ 4+48,\ 4+96,\ 4+192
\]
womit die Abstände übereinstimmen sollten, gar nicht einmal eine continuirliche Reihe ist. Das Glied, was vor $4+3$ hergeht, muss ja nicht $4$, das ist $4+0$, sondern $4+1\frac12$ sein. Also zwischen $4$ und $4+3$ sollten noch unendlich viele zwischen liegen, oder wie \textsc{Wurm} es ausdrückt (M.\ C.\ 1801 Juni St.\ S.\ 594) für Mercur oder für $n=1$ kommt aus $4+2^{n-2}\,3$ nicht $4$ sondern $5\frac12$.

„Es ist gar nicht zu tadeln, wenn man dergleichen ungefähre Übereinstimmungen in der Natur aufsucht. Die grössten Männer haben solchem \textit{lusus ingenii} nachgehängt. Aber so viel sich auch \textsc{Kepler} auf seine mit den Planeten-Distanzen in Übereinstimmung gebrachten regulären Körper zu gute that (er wollte, wie er sagte, die Ehre dieses Fundes nicht um das Churfürstenthum Sachsen geben): so hätte er doch gewiss den Planetismus des Uranus nicht damit angefochten (wenn diese Entdeckung zu seiner Zeit gemacht worden wäre), weil er nicht zu seinen Ideen passte. Er hätte vielmehr ohne allen Zweifel diese sogleich aufgegeben. Eben so hatten Sie es gemacht. Auch Sie haben über jenes Gesetz von \textsc{Titius} Betrachtungen angestellt; aber Sie haben sie nur Träume genannt, und sie sogleich fahren lassen, als sie anfingen, Thatsachen zu widersprechen.“
\begin{flushright}[\textsc{von Zach.}]\end{flushright}

\newcommand{\MercurySymbol}{\text{\mdseries\AstroSymbols ☿}}
\bigskip\Needspace{10\baselineskip}
\pdfbookmark[1]{Vorübergang des Mercur vor der Sonne, den 3. November 1802}{mercur18021103}
\begin{center}
VORÜBERGANG DES MERCUR VOR DER SONNE, DEN 3. NOVEMBER 1802.\par
\medskip\hrule\medskip
Monatliche Correspondenz. Band VI. S.\ 567—575. 1802 December.\par
\medskip\hrule
\end{center}\medskip
—\quad —\quad —

In \textit{Braunschweig} beobachtete Dr.\ \textsc{Gauss} die \textit{innere} Berührung beim Austritt um $0^{\mathrm u}\ 39^{\mathrm m}\ 16^{\mathrm s}$, die \textit{äussere} um $0^{\mathrm u}\ 40^{\mathrm m}\ 48^{\mathrm s}$ wahre Zeit. [Handschrift: Austritt des $\MercurySymbol$ beobachtet zu Braunschweig Nov.\ 9.\ $0^{\mathrm u}\ 23^{\mathrm m}\ 17^{\mathrm s}$ innere Berührung, $0^{\mathrm u}\ 24^{\mathrm m}\ 49^{\mathrm s}$ äussere Berührung Mittlere Zeit.] Er bediente sich dazu eines von \textsc{Baumann} verfertigten zweifüssigen Achromaten, in welchen Mercur und die schöne Gruppe von Sonnenflecken, an der er seinen Weg nahm, sehr gut ins Auge fiel. Aber leider war er nicht im Stande, eine so gute Zeitbestimmung zu machen, als er wohl gewünscht hatte. Diesen Unfall hat Dr.\ \textsc{Gauss}
\SourceLeaf{237}%
bei gegenwärtiger Jahreszeit mit mehreren Astronomen gemein, welche mit keinem Passagen-Instrumente versehen sind. Indessen musste er sich mit einigen, nach 1 Uhr, auf einem unbedeckten Öl Horizont genommenen einzelnen Sonnenhöhen zur Zeitbestimmung begnügen. Aus dieser Ursache, und auch deswegen, weil dies die erste Beobachtung dieser Art ist, die er gemacht hat, bittet er um Nachsicht. wenn seine Beobachtung mit anderen nicht genau übereinstimmen sollte. Ausser dem Austritte hat er auch das Fernrohr als Kreis-Mikrometer zu brauchen versucht, und mehrere Ascensional-Differenzen zwischen Mercur und dem Mittelpunkte der Sonne genommen, welchen er aber selbst keinen grossen Grad von Genauigkeit zutrauet. Indessen können sie immer durch Rechnung geprüft werden.
\begin{center}
\begin{tabular}{r|r}
Wahre Zeit&\shortstack{Diff. AR.\\$\odot-\MercurySymbol$}\\\hline
$23^{\mathrm u}\ 26^{\mathrm m}\ 6^{\mathrm s}$&$7'\ 51''$\\
$23\ 42\ 30$&$9\ 29$\\
$23\ 50\ 28$&$10\ 3$\\
$0\ 3\ 25$&$11\ 37$\\
$0\ 25\ 14$&$13\ 41$
\end{tabular}
\end{center}
—\quad —\quad —
\begin{flushright}[\textsc{von Zach.}]\end{flushright}

\bigskip\Needspace{10\baselineskip}
\pdfbookmark[1]{Pallas Olbersiana. 1802 December}{pallas180212}
\begin{center}
PALLAS OLBERSIANA.\par
\medskip\hrule\medskip
Monatliche Correspondenz. Band VI. S.\ 579—583. 1802 December.\par
\medskip\hrule
\end{center}\medskip
—\quad —\quad —

Diese sämmtlichen Mailänder Beobachtungen \textsc{Oriani}'s hat nun Dr.\ \textsc{Gauss} mit seinen III. Elementen der Bahn [S.\ 218 d.\ B.] dieses Planeten verglichen und nachstehende Übereinstimmung gefunden.
\Needspace{22\baselineskip}
\begin{center}\textit{Vergleichung der \textsc{Oriani}'schen Beobachtungen am Mailänder Aequatorial-Sector mit Dr.\ \textsc{Gauss}'s III. Elementen der Bahn.}\end{center}
\begin{center}
\begingroup\footnotesize\setlength{\tabcolsep}{2pt}
\begin{tabular}{l|r|r|r|r|r}
1802&\shortstack{Mittl. Zeit\\in Mailand}&\shortstack{Berechnete\\AR.\\der Pallas}&\shortstack{Unterschied\\R—B}&\shortstack{Berechnete\\nördl. Ab-\\weich. der $\PallasSymbol$}&\shortstack{Unterschied\\R—B}\\\hline
Juli 10&$9^{\mathrm u}\ 40^{\mathrm m}\ 23^{\mathrm s}$&$191^0\ 27'\ 33''$&$\left\{\begin{array}{r}-11''\\-35\end{array}\right.$&$17^0\ 56'\ 20''$&$\left\{\begin{array}{r}-4''\\+4\end{array}\right.$\\
18&$9\ 36\ 4$&$193\ 43\ 47$&$-17$&$17\ 2\ 29$&$-38$\\
24&$9\ 32\ 16$&$195\ 30\ 59$&$-29$&$16\ 19\ 59$&$+12$\\
25&$9\ 44\ 12$&$195\ 49\ 25$&$-23$&$16\ 12\ 40$&$+12$\\
28&$9\ 8\ 6$&$196\ 44\ 25$&$\left\{\begin{array}{r}-33\\-28\end{array}\right.$&$15\ 50\ 56$&$\left\{\begin{array}{r}+1\\+16\end{array}\right.$\\
29&$9\ 11\ 28$&$197\ 3\ 9$&$\left\{\begin{array}{r}-36\\-22\\-36\end{array}\right.$&$15\ 43\ 33$&$\left\{\begin{array}{r}+1\\+18\\+14\end{array}\right.$\\
31&$9\ 6\ 9$&$197\ 40\ 46$&$-34$&$15\ 28\ 55$&$-1$\\
Aug.\ 1&$9\ 4\ 14$&$197\ 59\ 44$&$\left\{\begin{array}{r}-40\\-45\\-49\end{array}\right.$&$15\ 21\ 18$&$\left\{\begin{array}{r}-20\\+2\\+5\end{array}\right.$\\
2&$9\ 23\ 37$&$198\ 19\ 4$&$\left\{\begin{array}{r}-36\\-41\end{array}\right.$&$15\ 13\ 43$&$\left\{\begin{array}{r}+10\\+12\end{array}\right.$
\end{tabular}\endgroup
\end{center}
\SourceLeaf{238}%
\begin{center}
\begingroup\footnotesize\setlength{\tabcolsep}{2pt}
\begin{tabular}{l|r|r|r|r|r}
1802&\shortstack{Mittl. Zeit\\in Mailand}&\shortstack{Berechnete\\AR.\\der Pallas}&\shortstack{Unterschied\\R—B}&\shortstack{Berechnete\\nördl. Ab-\\weich. der $\PallasSymbol$}&\shortstack{Unterschied\\R—B}\\\hline
Aug.\ 4&$8^{\mathrm u}\ 55^{\mathrm m}\ 38^{\mathrm s}$&$198\ 57\ 9$&$\left\{\begin{array}{r}-39''\\-47\end{array}\right.$&$14\ 58\ 52$&$\left\{\begin{array}{r}-5''\\-4\end{array}\right.$\\
5&$9\ 14\ 57$&$199\ 16\ 47$&$-39$&$14\ 51\ 13$&$\left\{\begin{array}{r}+15\\+28\end{array}\right.$\\
6&$9\ 6\ 27$&$199\ 36\ 7$&$\left\{\begin{array}{r}-50\\-51\end{array}\right.$&$14\ 43\ 43$&$-3$\\
7&$9\ 4\ 18$&$199\ 55\ 38$&$\left\{\begin{array}{r}-48\\-47\end{array}\right.$&$14\ 36\ 8$&$+15$\\
8&$9\ 0\ 14$&$200\ 15\ 9$&$-56$&$14\ 28\ 35$&$+14$
\end{tabular}\endgroup
\end{center}

Dr.\ \textsc{Gauss} hat zwar späterhin (October-St.\ S.\ 394.) [S.\ 223 d.\ B.] zum IV. mal verbesserte Elemente gegeben, aber obige Mailänder Beobachtungen nicht selbst damit verglichen; indessen hat er doch folgende Vergleichung angestellt, aus welcher sich beiläufig beurtheilen lässt, in wie fern diese IV. Elemente besser mit den Beobachtungen stimmen, als die III.
\begin{center}
Die IV. Elemente geben\par\medskip
\begin{tabular}{l|r|r}
1802&Die AR. grösser&Die Declination kleiner\\\hline
Juli 8&$24''$&$8''$\\
Aug.\ 4&$66$&$0$
\end{tabular}\par\medskip
als die III. Elemente.
\end{center}

So wie also die III. Elemente um die Zeit des Schlusses der Beobachtungen die AR. gegen 1 Minute zu klein geben, so geben die IV. Elemente dieselbe etwa $20''$ zu gross. Die Declinationen stimmen bei beiden, so weit es die Genauigkeit der Beobachtungen zulässt, ziemlich gut.

Dr.\ \textsc{Gauss} hat nun auch nach diesen letztern \textsc{Oriani}'schen Beobachtungen die Elemente dieser Bahn zum V. mal verbessert, und nachstehende Bestandtheile erhalten.
\Needspace{14\baselineskip}
\begin{center}\textit{V. Elemente der Pallas Olbersiana.}\end{center}
\begin{center}
\begingroup\small\setlength{\tabcolsep}{3pt}
\begin{tabular}{l r}
\shortstack[l]{Epoche 1802 März 31 Mittag\\in Seeberg}&$162^0\ 46'\ 58''2$\\
Sonnenferne $\ .\ .\ .\ .\ .\ .\ .\ .\ .$&$301\ 28\ 24.0$\\
Aufsteigenden Knoten $\ .\ .\ .\ .\ .\ .$&$172\ 27\ 3.0$\\
Tägliche mittl. tropische Bewegung $\ .$&$769''583$\\[3pt]
Logarithmus der halben grossen Axe&$0.4425529$\\
Excentricität $\ .\ .\ .\ .\ .\ .\ .\ .\ .$&$0.244976$\\
Neigung der Bahn $\ .\ .\ .\ .\ .\ .\ .$&$34^0\ 37'\ 40''$
\end{tabular}\endgroup
\end{center}

Mit diesen Elementen verglich nun Dr.\ \textsc{Gauss} die späteste aller Beobachtungen von \textsc{Messier} (November-St.\ S.\ 502.) und fand:
\Needspace{12\baselineskip}
\begin{center}
\begingroup\small\setlength{\tabcolsep}{4pt}
\begin{tabular}{l|r|r|r}
1802&\shortstack{Mittlere Zeit\\in Paris}&\shortstack{Berechnete AR.\\der Pallas}&\shortstack{Unterschied\\R—B}\\\hline
Sept.\ 21&$7^{\mathrm u}\ 28^{\mathrm m}$&$215^0\ 49'\ 14''$&$+28''$
\end{tabular}\par\medskip
\begin{tabular}{r|r|r|r}
\shortstack{Berechnete Declination\\der Pallas}&\shortstack{Unterschied\\R—B}&\shortstack{Abstand\\von der $\EarthSymbol$}&Lichtstärke\\\hline
$9^0\ 0'\ 16''\ \mathrm N.$&$+48''$&$3.518$&$0.485$
\end{tabular}\endgroup
\end{center}

Der Unterschied der Rechnung und Beobachtung ist also noch mässig. Dr.\ \textsc{Gauss} schreibt daher:
„Da diese \textsc{Messier}'sche Beobachtung in einer so wenig günstigen Lage gemacht worden ist (nach einer beiläufigen Rechnung hatte Pallas nur 13 Grad Höhe über dem Horizont), so bin ich ganz Ihrer Meinung, dass man von dieser \textit{einzelnen} Beobachtung nicht viel sicheres zur Verbesserung der Elemente hernehmen kann. Inzwischen würde die Verbesserung auf den Ort der \textit{Pallas} im künftigen Jahre wol eben keinen sehr bedeutenden Einfluss haben. Ich habe mehrere Örter nach den V. Elementen neu
\SourceLeaf{239}%
berechnet, und die AR. für den 4. Febr., 24 März, 11. Mai, 28. Juni 1803 um $3\frac12$ Min., $5\frac12$ Min., $7\frac12$ Min., 9 Min. kleiner, die Abweichungen aber alle um ungefähr 1 Min. grösser gefunden, als in der nach den IV. Elementen berechneten Ephemeride.“ (October-St.\ S.\ 395.) [S.\ 224 d.\ B.]

Inzwischen um auch hier keinen Wunsch übrig zu lassen, so hat Dr.\ \textsc{Gauss} diese ganze Ephemeride des Laufes der Pallas im Jahr 1803 von neuem nach diesen V. Elementen berechnet, auch noch auf einen Monat weiter ausgedehnt. Da \textit{Pallas} am 3. April 1803 anfängt, eine grössere Lichtstärke zu erhalten, als sie am 21. September dieses Jahres hatte, und man sie sodann in einer viel grösseren Höhe über dem Horizont beobachten kann, so darf man jetzt wol nicht mehr daran zweifeln, dass sie sich unseren stärkeren Fernröhren nicht entziehen wird. Hier folgt demnach: [Die Ephemeride von Febr.\ 4 bis Aug.\ 9. 1803, welche AR und Decl. so wie Lichtstärke in der bisher angewandten Form enthalten.]

In der dritten Columne ist zur Einheit diejenige Lichtstärke angenommen, welche der Planet in der Entfernung 1 von der Sonne und Erde haben würde. Nach diesem Massstabe war sie 1802
\begin{center}
\begingroup\small
\begin{tabular}{llr r l}
1802&April&4 $\ .\ .$&$0.08997$&[Handschrift:] Genauer $0.09019$\\
&Mai&16 $\ .\ .$&$0.04740$&\\
&Aug.&10 $\ .\ .$&$0.01455$&\\
&Sept.&21 $\ .\ .$&$0.01030$&
\end{tabular}\par\medskip
\begin{tabular}{llr r}
1801&Jan.&1 $\ .\ .$&$0.03637$\\
&Febr.&11 $\ .\ .$&$0.02409$\\
&Dec.&7 $\ .\ .$&$0.02494$\\
1802&Febr.&26 $\ .\ .$&$0.05633$
\end{tabular}\endgroup
\end{center}
\begin{flushright}[\textsc{von Zach.}]\end{flushright}

\bigskip\Needspace{12\baselineskip}
\pdfbookmark[1]{Gauss an Bode. 1803 März 3}{bode18030303}
\begin{center}
GAUSS AN BODE.\par
\medskip\hrule\medskip
Astronomisches Jahrbuch für 1806. S.\ 179—180. Berlin 1803.\par
\medskip\hrule
\end{center}\medskip
\begin{flushright}Braunschweig 1803 März 3.\end{flushright}

Die wichtige Nachricht von der so frühen Wiederauffindung des \textsc{Olbers}'schen Planeten zu Lilienthal werden Ew. — vermuthlich schon von Hrn.\ \textsc{Harding} selbst erhalten haben. So werden Sie wahrscheinlich auch schon im Besitz von Dr.\ \textsc{Olbers} Beobachtungen sein. Ich schmeichle mir, dass es Ihnen nicht unlieb sein wird, auch die Resultate davon näher zu erfahren. Das Wetter ist seit dem 23. Febr. hier sehr schlecht gewesen; der Mondschein kommt jetzt dazu. Die Beobachtungen werden also wol noch etwas unterbrochen bleiben.

Hier ist die Vergleichung der \textsc{Olbers}'schen Beobachtung mit meinen V. Elementen, wonach die Ephemeride im Decemberheft der M.\ C. berechnet worden ist.
\begin{center}
\begingroup\small\setlength{\tabcolsep}{3pt}
\begin{tabular}{l|r|r|r|r|r}
& &\multicolumn{2}{c|}{Berechnete}&\multicolumn{2}{c}{Unterschied R—B}\\
1803&\shortstack{Mittl. Zeit\\in Bremen}&AR.&Declin.&AR.&Declin.\\\hline
Febr.\ 21&$17^{\mathrm u}\ 6^{\mathrm m}\ 10^{\mathrm s}$&$272^0\ 58'\ 47''$&$7^0\ 30'\ 40''$&$+2'\ 2''$&$-34''$\\
23&$15\ 24\ 36$&$273\ 31\ 14$&$7\ 45\ 4$&$+2\ 35$&$-57$
\end{tabular}\endgroup
\end{center}

Ich habe meiner Begierde nicht widerstehen können, hienach eine vorläufige Verbesserung der Elemente vorzunehmen; ich behalte mir vor bessere zu geben, sobald zahlreichere und schärfere Beobachtungen vorhanden sein werden. Ew. — werden aus der Vergleichung mit den V. [S.\ 233 d.\ B.] sehen dass der Unterschied nur sehr gering ist:
\SourceLeaf{240}%
\begin{center}
\begingroup\small\setlength{\tabcolsep}{3pt}
\begin{tabular}{l r}
Epoche 1803 in Seeberg. Meridian $\ .\ .$&$221^0\ 28'\ 54''0$\\
Jährliche trop. Bewegung (365 T.) $\ .\ .$&$78\ 0\ 36.9$\\
Tägliche $\ .\ .\ .\ .\ .\ .\ .\ .\ .\ .\ .$&$769''416$\\
Sonnenferne 1803 siderisch ruhend $\ .$&$301^0\ 24'\ 13''$\\[3pt]
Knoten $\ .\ .\ .\ .\ .\ .\ .\ .\ .\ .\ .$&$172^0\ 28'\ 8''$\\
Excentricität $\ .\ .\ .\ .\ .\ .\ .\ .\ .$&$0.245619$\\
Logarithmus der halben grossen Axe&$0.4426160$\\
Neigung der Bahn $\ .\ .\ .\ .\ .\ .\ .\ .$&$34^0\ 38'\ 20''$
\end{tabular}\endgroup
\end{center}

Diese Elemente stellen die Beobachtung folgendermassen dar:
\begin{center}
\begingroup\small\setlength{\tabcolsep}{4pt}
\begin{tabular}{l|r|r|r|r}
&\multicolumn{2}{c|}{Berechnete}&\multicolumn{2}{c}{Unterschied R—B}\\
1803&AR.&Declin.&AR.&Declin.\\\hline
Febr.\ 21&$272^0\ 56'\ 25''0$&$7^0\ 31'\ 29''8$&$-20''0$&$+15''8$\\
23&$273\ 28\ 46.8$&$7\ 45\ 52.3$&$+7.8$&$-7.7$
\end{tabular}\endgroup
\end{center}
Man wird demnach den Planeten in diesem Jahre immer mit aller zu verlangenden Leichtigkeit finden können.

Hoffentlich wird auch die $\CeresSymbol$ nun bald aufgefunden werden. Die Auffindung derselben wird diesmal wegen ihres niedrigen Standes am südlichen Himmel wol in Italien und Frankreich früher gelingen als bei uns.

\bigskip\Needspace{13\baselineskip}
\pdfbookmark[1]{Tafeln für die Störungen der Ceres. 1803 März}{ceresstoerungen180303}
\begin{center}
TAFELN FÜR DIE STÖRUNGEN DER CERES.\par
\medskip\hrule\medskip
Monatliche Correspondenz. Band VII. S.\ 259—262. 1803 März.\par
\medskip\hrule\bigskip
GAUSS AN VON ZACH.
\end{center}\medskip
\begin{flushright}Braunschweig, den 26. December 1802.\end{flushright}

Vor einiger Zeit schrieb ich Ihnen. dass ich auf eine Abkürzung der Tafeln für die Störungen der \textit{Ceres} gedacht habe und Ihnen davon eine Probe künftig vorlegen wolle. (Mon. Corr. VI. B. S.\ 498) [230 d.\ B.] Ich habe jetzt die Ehre, mich darüber näher zu erklären, und Ihnen nicht eine Probe, sondern eine Tafel \textit{in extenso} zu übersenden.

Die Anzahl der Gleichungen ist schon jetzt, da ich mich noch auf die Störungen durch Jupiter und auf die ersten Potenzen der Excentricitäten eingeschränkt habe, so ansehnlich, dass es äusserst beschwerlich sein würde, ihre numerischen Werthe für eine grosse Anzahl von Örtern ohne Tafeln aus den Formeln selbst zu berechnen. Aber auch mit Tafeln, nach der gewöhnlichen Art eingerichtet, ist die Arbeit wahrlich beschwerlich. Für \textit{alle} Gleichungen, die im November-Heft 1802 der M.\ C. [228 d.\ B.] abgedruckt sind, würden \textit{vierzig} Tafeln nöthig sein, und noch fast \textit{dreissig}, wenn man die kleinern Gleichungen, die einzeln unter $2''$ betragen, weglässt. So viele Argumente zu bilden, in so viele Tafeln einzugehen, und so viele Additionen zu machen, erfordert noch immer viel Geduld, wenn es oft kommt
Ich habe daher diese Tafeln auf eine geschmeidigere Art abzukürzen gesucht. Die \textit{acht} ersten Gleichungen \textit{für die Länge} (wo beiläufig gesagt, der erste Coëfficient S.\ 495 [228] anstatt $230''00$, $231''00$ sein muss), kommen von selbst in \textit{eine} Tafel, deren Argument $\CeresSymbol-\JupiterSymbol$ ist. Aber die vierzehn übrigen würden eben so viele Tafeln erfordern. Man könnte sie in eine Tafel mit doppelten Eingängen bringen: die würde aber so weitläuftig, mühsam zu berechnen und beschwerlich zu gebrauchen sein, dass daran nicht zu denken ist. Ich habe daher folgendes Verfahren gewählt:
\SourceLeaf{241}%

Bezeichne ich Kürze halber $\CeresSymbol-\JupiterSymbol$ mit $D$, so sieht man bei einiger Aufmerksamkeit, dass jede dieser vierzehn Gleichungen sich unter die Form
\[
a\sin(iD-\CeresSymbol-A)=a\sin(iD-A)\cos\CeresSymbol-a\cos(iD-A)\sin\CeresSymbol
\]
bringen lässt, wo $i$ positive und negative ganze Zahlen vorstellt. Die Summe aller Gleichungen wird also durch $P\cos\CeresSymbol-Q\sin\CeresSymbol$ ausgedrückt werden, wo $P$ und $Q$ Functionen von $D$ sind. Macht man
\[
\frac PQ=\operatorname{tang}\varphi,\quad
\frac P{\sin\varphi}=\frac Q{\cos\varphi}=R,
\]
so sind also auch $R$ und $\varphi$ blos von $D$ abhängig, und die Summe aller vierzehn Gleichungen
\[
=R\sin(\varphi-\CeresSymbol)
\]
Man kann also in einer Tafel, deren Argument $D$ ist, für alle Werthe von $D$ die entsprechenden Werthe von $R$ und $\varphi$ ansetzen, und findet dadurch vermittelst einer sehr leichten Rechnung sogleich den Inbegriff aller jener Gleichungen. Ich habe indess vorgezogen, nicht $\varphi$ sondern $\varphi-2D$ anzusetzen, weil dieser Winkel sich nicht so stark ändert als $\varphi$, und innerhalb gewisser Gränzen nur periodische Oscillationen macht, während dass $D$ alle Werthe von $0$ bis $360^0$ durchläuft, da hingegen $\varphi$ unterdess zwei Umläufe machen würde. Dieses $\varphi-2D$ ist das, was in der dritten Tafel mit $B$ bezeichnet ist, so wie obiges $R$ mit $A$. Auf diese Weise ist also der Inbegriff aller von der einfachen Excentricität abhängigen Gleichungen der Länge
\[
A\sin(B+2D-\CeresSymbol)=A\sin(B+\CeresSymbol-2\JupiterSymbol)
\]

Auf ähnliche Weise sind die von der Excentricität unabhängigen Gleichungen des Radius Vector in der Tafel II enthalten, und die von den einfachen Excentricitäten abhängigen in der Tafel IV unter die Form $A\sin(B+2\CeresSymbol-3\JupiterSymbol)$ gebracht; endlich unter dieselbe Form in Tafel V alle im November-Heft [S.\ 229 d.\ B.] aufgeführten Gleichungen der Breite.

Auf diese Art haben wir also nur fünf Tafeln, wobei freilich etwas trigonometrische Rechnung erforderlich ist, die aber mir wenigstens weit bequemer fällt, als wenn ich die einzelnen hier zusammengefassten Gleichungen alle aus besondern Tafeln nehmen müsste. Auch liessen sich dabei mit Vortheil noch einige Abänderungen treffen, wenn man z.\ B. statt $A$ gleich $\log A$ ansetzte, \textit{oder auch auf andere Weise}. Dieses will ich mir aber auf die Zukunft ersparen, da überhaupt diese Tafeln nur so lange dienen sollen, bis mit künftigen Beobachtungen die Elemente noch ferner berichtigt und die Störungen danach genauer und vollständiger berechnet sein werden. Vor der Hand denke ich, dass mir diese Tafeln, besonders bei der Berechnung der Beobachtungen des künftigen Jahres gute Dienste leisten sollen.

\Needspace{7\baselineskip}
\begin{center}\textit{Anmerkungen.}\end{center}
1.\quad $\CeresSymbol$ und $\JupiterSymbol$ bedeuten die mittlern heliocentrischen Längen der Ceres und des Jupiter.

2.\quad Die Störungen des Radius Vector sind in Million-Theilen des Halb\-messers der Erdbahn ausgedrückt.

3.\quad Bei den Störungen der Breite, und bei den zweiten Theilen der Störungen der Länge und des Radius Vector ist $A$ als positiv anzusehen, und bei dem Sinus von $B+\CeresSymbol-2\JupiterSymbol$ und $B+2\CeresSymbol-3\JupiterSymbol$ das Zeichen gehörig zu beobachten.

4.\quad Bei der Breite nimmt man $B$ aus der ersten Columne für die ersten sechs Zeichen des Arguments oder beim Niedersteigen; aus der dritten Columne hingegen in den sechs letzten Zeichen oder beim Aufsteigen; $A$ ist beiden gemeinschaftlich. Die Anordnung der Tafel zeigt dies schon von selbst.
